\documentclass[12pt,oneside]{book}

%============================================================
% PACKAGES
%============================================================

%--------------------
% Page + font
%--------------------
\usepackage[margin=1in]{geometry}
\usepackage[T1]{fontenc}
\usepackage[utf8]{inputenc}
\usepackage{lmodern}
\usepackage{microtype}
\usepackage{parskip}

%--------------------
% Math
%--------------------
\usepackage{amsmath, amssymb, amsfonts, amsthm, mathtools}
\usepackage{bm}
\usepackage{bbm}
\usepackage{mathrsfs}
\usepackage{dsfont}

%--------------------
% Tables + figures
%--------------------
\usepackage{graphicx}
\usepackage{float}
\usepackage{booktabs}
\usepackage{array}
\usepackage{xltabular}
\usepackage{multirow}
\usepackage{caption}
\usepackage{subcaption}

%--------------------
% Lists
%--------------------
\usepackage{enumitem}
\setlist[itemize]{topsep=2pt,itemsep=2pt}
\setlist[enumerate]{topsep=2pt,itemsep=2pt}

%--------------------
% Colors + hyperlinks
%--------------------
\usepackage[dvipsnames]{xcolor}
\usepackage[
    colorlinks=true,
    linkcolor=MidnightBlue,
    citecolor=MidnightBlue,
    urlcolor=MidnightBlue
]{hyperref}
\usepackage[nameinlink,capitalise]{cleveref}

%--------------------
% Algorithms
%--------------------
\usepackage{algorithm}
\usepackage{algpseudocode}

%--------------------
% Code blocks
%--------------------
\usepackage{listings}

\lstdefinestyle{codestyle}{
    basicstyle=\ttfamily\small,
    breaklines=true,
    frame=single,
    numbers=left,
    numberstyle=\tiny,
    keywordstyle=\color{MidnightBlue},
    commentstyle=\color{ForestGreen},
    stringstyle=\color{BrickRed},
    showstringspaces=false,
    tabsize=4
}

\lstset{style=codestyle}

%--------------------
% Section styling
%--------------------
\usepackage{titlesec}
\usepackage{tocloft}

\setcounter{tocdepth}{2}
\setcounter{secnumdepth}{3}

%============================================================
% THEOREMS
%============================================================

\numberwithin{equation}{section}

\theoremstyle{definition}
\newtheorem{definition}{Definition}[section]
\newtheorem{example}{Example}[section]
\newtheorem{exercise}{Exercise}[section]
\newtheorem{remark}{Remark}[section]

\theoremstyle{plain}
\newtheorem{theorem}{Theorem}[section]
\newtheorem{lemma}{Lemma}[section]
\newtheorem{proposition}{Proposition}[section]
\newtheorem{corollary}{Corollary}[section]

\theoremstyle{remark}
\newtheorem*{note}{Note}

%============================================================
% CUSTOM COMMANDS
%============================================================

% Sets
\newcommand{\R}{\mathbb{R}}
\newcommand{\N}{\mathbb{N}}
\newcommand{\Z}{\mathbb{Z}}
\newcommand{\Q}{\mathbb{Q}}
\newcommand{\C}{\mathbb{C}}

% Probability
\newcommand{\E}{\mathbb{E}}
\newcommand{\Var}{\mathrm{Var}}
\newcommand{\Cov}{\mathrm{Cov}}
\newcommand{\Corr}{\mathrm{Corr}}
\newcommand{\Prob}{\mathbb{P}}

\newcommand{\ind}{\mathbbm{1}}
\newcommand{\iid}{\overset{\text{iid}}{\sim}}
\newcommand{\convd}{\overset{d}{\to}}
\newcommand{\convp}{\overset{p}{\to}}
\newcommand{\as}{\overset{a.s.}{\to}}

% Operators
\newcommand{\given}{\,\middle|\,}
\newcommand{\argmin}{\operatorname*{arg\,min}}
\newcommand{\argmax}{\operatorname*{arg\,max}}

% Linear algebra / analysis
\newcommand{\norm}[1]{\left\lVert #1 \right\rVert}
\newcommand{\abs}[1]{\left\lvert #1 \right\rvert}
\newcommand{\inner}[2]{\left\langle #1, #2 \right\rangle}

\newcommand{\dd}{\,\mathrm{d}}
\newcommand{\grad}{\nabla}
\newcommand{\Hess}{\nabla^2}

\newcommand{\tr}{\operatorname{tr}}
\newcommand{\rank}{\operatorname{rank}}
\newcommand{\diag}{\operatorname{diag}}

% Distributions
\newcommand{\Normal}{\mathcal{N}}
\newcommand{\Bern}{\mathrm{Bernoulli}}
\newcommand{\Bin}{\mathrm{Binomial}}
\newcommand{\Pois}{\mathrm{Poisson}}
\newcommand{\Gam}{\mathrm{Gamma}}
\newcommand{\BetaDist}{\mathrm{Beta}}
\newcommand{\Exp}{\mathrm{Exponential}}
\newcommand{\Dir}{\mathrm{Dirichlet}}
\newcommand{\InvGam}{\mathrm{Inv\text{-}Gamma}}

% Title block
\title{EN.553.781: Functional Data Analysis}
\author{Jonathan Y. Ma}
\date{June 2026}

\begin{document}

\maketitle

\noindent \emph{This course was taught by Dr. Anuj Srivastava in the Fall 2026 Semester. The notes are transcribed by Jonathan Ma and may contain errors and omitted content, so any issues are to be directed towards me.}

\tableofcontents

\section{Functional and Shape Data}

\begin{definition}[Functional datum]
A functional datum is an observation whose basic statistical unit is a
function
\[
    f:D\longrightarrow\R^q,
\]
rather than a finite list of unrelated measurements. The domain $D$ is often
an interval of time, but may also be a spatial region or a manifold.
\end{definition}

A sample in functional data analysis (FDA) is therefore
$f_1,\ldots,f_n$, where each $f_i$ is viewed as one observation. In practice
we observe noisy, discrete samples
\[
    Y_{ij}=f_i(t_{ij})+\varepsilon_{ij},
\]
possibly at subject-specific times $t_{ij}$. Smoothing and registration turn
these measurements into objects on a common functional domain.

FDA supports familiar statistical tasks:
\begin{itemize}
    \item means, medians, covariance, and principal modes of variation;
    \item signal estimation from noisy measurements;
    \item testing, classification, clustering, and ANOVA;
    \item regression with functional predictors, responses, or both.
\end{itemize}
The distinction from multivariate statistics is not merely that the vectors
are long. Functions live in infinite-dimensional spaces and carry geometric
structure such as smoothness, derivatives, and time transformations.

\subsection{FDA Versus Multivariate and Time-Series Analysis}

If every curve is observed on the same reliable grid
$t_1,\ldots,t_m$, it can be stored as
\[
    (f_i(t_1),\ldots,f_i(t_m))^\mathsf T\in\R^m.
\]
This representation permits ordinary multivariate tools, but its results may
depend on grid density and ignore smoothness. A functional representation
allows interpolation, differentiation, integration, irregular sampling, and
comparison across grids.

Time-series models typically emphasize sequential dynamics, for example
$X_{j+1}=h(X_j)+\varepsilon_j$, with ordered and often equally spaced times.
FDA typically treats an entire trajectory as one observation and models
cross-sectional variation among trajectories. The two perspectives overlap,
but answer different questions.

\begin{example}[Growth curves]
For child $i$, height measurements $Y_{ij}$ are taken at ages $t_{ij}$. The
scientific objects may be the smooth growth curve $f_i(t)$, its velocity
$f_i'(t)$, and population variation in the timing of growth spurts. Treating
the measurements as an unordered vector loses this derivative and timing
structure.
\end{example}

\subsection{Shape Data and Invariance}

Shape analysis is closely related to FDA because curves and surfaces are
represented by functions. Following Kendall's principle, shape is what
remains after removing transformations that should not change the object,
commonly translation, rotation, and scale.

\begin{definition}[Group invariance]
Let a group $G$ act on a representation space $\mathcal F$. A quantity
$S:\mathcal F\to\mathcal S$ is invariant to $G$ if
\[
    S(g\cdot f)=S(f)\qquad\text{for every }g\in G.
\]
Two objects represent the same shape when they lie in the same orbit
$[f]=\{g\cdot f:g\in G\}$.
\end{definition}

A shape distance should be defined on equivalence classes, not arbitrary
representatives. This motivates quotient spaces later in the course. Typical
shape tasks include distances, optimal deformations or geodesics, means,
principal modes, and population comparisons.

\section{Metric and Vector Spaces}

\begin{definition}[Metric space]
A metric on a set $\mathcal X$ is a function
$d:\mathcal X\times\mathcal X\to[0,\infty)$ such that for all $x,y,z$:
\begin{enumerate}
    \item $d(x,y)=0$ if and only if $x=y$;
    \item $d(x,y)=d(y,x)$;
    \item $d(x,z)\leq d(x,y)+d(y,z)$.
\end{enumerate}
The pair $(\mathcal X,d)$ is a metric space.
\end{definition}

Examples include Euclidean distance on $\R^p$ and the uniform distance
\[
    d_\infty(f,g)=\sup_{t\in[0,1]}\abs{f(t)-g(t)}
\]
on $C([0,1])$. A metric determines convergence, continuity, neighborhoods,
and consequently what it means for two functional observations to be close.

\begin{definition}[Vector space]
A real vector space $V$ is a set closed under addition and scalar
multiplication satisfying associativity, commutativity of addition,
distributivity, a zero element, and additive inverses.
\end{definition}

Both $\R^p$ and $C([0,1])$ are vector spaces under pointwise operations. If
$f,g\in C([0,1])$, then
$(af+bg)(t)=af(t)+bg(t)$ is again continuous.

\begin{definition}[Linear independence and basis]
Vectors $v_1,\ldots,v_m$ are linearly independent if
\[
    \sum_{j=1}^ma_jv_j=0\quad\Longrightarrow\quad a_1=\cdots=a_m=0.
\]
A basis is a linearly independent spanning set.
\end{definition}

Finite-dimensional spaces have finite bases. Function spaces generally
require infinite systems, and convergence of an infinite expansion must be
specified using a norm.

\begin{definition}[Linear transformation]
For vector spaces $V,W$, a map $T:V\to W$ is linear when
\[
    T(au+bv)=aT(u)+bT(v)
\]
for every $u,v\in V$ and scalars $a,b$.
\end{definition}

Differentiation, integration, and multiplication by a fixed matrix are linear.
The map $f\mapsto f^2$ is not.

\section{Normed and Banach Spaces}

\begin{definition}[Norm]
A norm on $V$ is a map $\norm{\cdot}:V\to[0,\infty)$ satisfying
\[
 \norm v=0\Longleftrightarrow v=0,\qquad
 \norm{av}=\abs a\norm v,\qquad
 \norm{u+v}\leq\norm u+\norm v.
\]
It induces the metric $d(u,v)=\norm{u-v}$.
\end{definition}

Important functional norms on $[0,1]$ are
\[
 \norm f_\infty=\sup_t\abs{f(t)},\qquad
 \norm f_p=\left(\int_0^1\abs{f(t)}^p\dd t\right)^{1/p}.
\]
The $L^2$ norm measures integrated squared size, while the sup norm controls
the worst pointwise deviation. Convergence in one need not imply convergence
in the other.

\begin{definition}[Cauchy sequence and completeness]
A sequence $\{x_n\}$ is Cauchy if for every $\varepsilon>0$ there is $N$ such
that $d(x_m,x_n)<\varepsilon$ whenever $m,n\geq N$. A complete metric space
contains the limit of every Cauchy sequence. A complete normed vector space is
called a Banach space.
\end{definition}

The space $C([0,1])$ with $\norm{\cdot}_\infty$ is Banach. The space
$L^p([0,1])$, formed by identifying functions equal almost everywhere, is
Banach for $1\leq p\leq\infty$.

\begin{remark}
The set of continuous functions equipped with the $L^2$ norm is not complete:
an $L^2$-Cauchy sequence of continuous functions may converge to a
discontinuous equivalence class. Its completion is $L^2([0,1])$.
\end{remark}

\section{Inner-Product and Hilbert Spaces}

\begin{definition}[Inner product]
An inner product on a real vector space $V$ is a map
$\inner{\cdot}{\cdot}:V\times V\to\R$ that is linear in its first argument,
symmetric, and positive definite:
\[
    \inner v v\geq0,\qquad \inner v v=0\Longleftrightarrow v=0.
\]
It induces $\norm v=\sqrt{\inner v v}$.
\end{definition}

On $\R^p$, $\inner x y=x^\mathsf Ty$. On $L^2([0,1])$,
\begin{equation}
    \inner f g_{L^2}=\int_0^1f(t)g(t)\dd t.
    \label{eq:l2-inner-product}
\end{equation}

\begin{theorem}[Cauchy--Schwarz]
For every inner-product space,
\[
    \abs{\inner u v}\leq\norm u\norm v.
\]
\end{theorem}

\begin{corollary}[Triangle inequality]
The norm induced by an inner product satisfies
$\norm{u+v}\leq\norm u+\norm v$.
\end{corollary}

\begin{definition}[Hilbert space]
A Hilbert space is a complete inner-product space.
\end{definition}

The space $L^2([0,1])$ is the canonical Hilbert space for elementary FDA. Its
geometry supports angles, orthogonality, projections, and least squares just
as Euclidean space does.

\subsection{Orthogonality and Projection}

Elements $u,v$ are orthogonal, written $u\perp v$, if $\inner u v=0$. An
orthonormal set $\{e_j\}$ satisfies
$\inner{e_j}{e_k}=\ind\{j=k\}$.

\begin{theorem}[Projection theorem]
Let $M$ be a closed linear subspace of a Hilbert space $H$. For every $f\in H$
there is a unique $m^*\in M$ minimizing $\norm{f-m}$. It is characterized by
\[
    f-m^*\perp M.
\]
\end{theorem}

For $M=\operatorname{span}(e_1,\ldots,e_J)$ with an orthonormal basis,
\[
    P_Mf=\sum_{j=1}^J\inner f{e_j}e_j.
\]
This is the infinite-dimensional origin of the normal equations in least
squares.

\subsection{Orthonormal Bases and Fourier Expansion}

\begin{definition}[Complete orthonormal system]
An orthonormal sequence $\{e_j\}_{j\geq1}$ is complete if the only element
orthogonal to every $e_j$ is zero. Equivalently, its linear span is dense in
the Hilbert space.
\end{definition}

\begin{theorem}[Hilbert-space expansion]
For a complete orthonormal system and $f\in H$,
\[
 f=\sum_{j=1}^\infty\inner f{e_j}e_j
 \quad\text{in }H,\qquad
 \norm f^2=\sum_{j=1}^\infty\abs{\inner f{e_j}}^2.
\]
\end{theorem}

The second identity is Parseval's identity. On $L^2([0,1])$, a real Fourier
system consists of the constant function and appropriately normalized
$\sin(2\pi kt)$ and $\cos(2\pi kt)$. Fourier bases are natural for periodic
functions; splines are usually more local and flexible near boundaries.

\begin{example}[Truncation]
The approximation
\[
    f_J(t)=\sum_{j=1}^Jc_je_j(t),\qquad c_j=\inner f{e_j},
\]
is the best $J$-term approximation within
$\operatorname{span}(e_1,\ldots,e_J)$ under $L^2$ loss. Smooth functions
typically have rapidly decaying coefficients in a compatible basis.
\end{example}

\section{From Discrete Measurements to Smooth Functions}

For one subject, suppose $(t_i,y_i)$ satisfy
$y_i=f(t_i)+\varepsilon_i$. Choose basis functions
$b_1,\ldots,b_J$ and approximate
\[
    f(t)=\sum_{j=1}^Jc_jb_j(t).
\]
Let $B_{ij}=b_j(t_i)$, $c=(c_1,\ldots,c_J)^\mathsf T$, and
$y=(y_1,\ldots,y_m)^\mathsf T$. Least-squares fitting becomes
\begin{equation}
    \widehat c=\argmin_c\norm{y-Bc}^2.
    \label{eq:fda-basis-ls}
\end{equation}
If $B$ has full column rank,
$\widehat c=(B^\mathsf TB)^{-1}B^\mathsf Ty$. In computation one should use a
QR or SVD solve rather than explicitly form the inverse.

The number and type of basis functions control the bias--variance tradeoff.
Too few cannot represent the signal; too many follow measurement noise. A
common coefficient vector has meaning only when all curves use the same basis
and domain.

\subsection{Roughness-Penalized Estimation}

Instead of choosing a very small basis, one may directly penalize roughness:
\begin{equation}
 \widehat f=\argmin_f
 \left\{\sum_{i=1}^m\{y_i-f(t_i)\}^2
 +\lambda\int_0^1\{f^{(r)}(t)\}^2\dd t\right\}.
 \label{eq:fda-roughness}
\end{equation}
Usually $r=2$, so curvature is penalized. The tuning parameter $\lambda$
controls the tradeoff:
\[
 \lambda\downarrow0:\text{ flexible fit},\qquad
 \lambda\uparrow\infty:\text{ approach the null space of }D^r.
\]
For $r=2$, the null space consists of linear functions.

With $f(t)=b(t)^\mathsf Tc$, define
\[
    R_{jk}=\int_0^1b_j^{(r)}(t)b_k^{(r)}(t)\dd t.
\]
Then
\begin{equation}
    \widehat c_\lambda
    =(B^\mathsf TB+\lambda R)^{-1}B^\mathsf Ty.
    \label{eq:penalized-basis-solution}
\end{equation}
This resembles ridge regression, but $R$ penalizes rough directions rather
than all coefficient directions equally.

\begin{proposition}[Linear-smoother form]
At the observation times, the penalized fitted values are
\[
    \widehat y=S_\lambda y,\qquad
    S_\lambda=B(B^\mathsf TB+\lambda R)^{-1}B^\mathsf T.
\]
Their effective degrees of freedom are
$\operatorname{df}(\lambda)=\tr(S_\lambda)$.
\end{proposition}

Cross-validation or generalized cross-validation can select $\lambda$. With
irregular or sparse subject-specific sampling, smoothing choices can strongly
affect later covariance estimates and should be propagated or at least
reported.

\begin{theorem}[Smoothing-spline representation]
For squared error with the integrated squared second-derivative penalty, the
minimizer over the appropriate Sobolev space is a natural cubic spline with
knots at the distinct observation times.
\end{theorem}

Thus an infinite-dimensional optimization reduces to a finite-dimensional
one without first imposing an arbitrary dense knot grid.

\section{Statistical Summaries of Functional Data}

Let $F$ be a random element of $L^2([0,1])$ with
$\E\norm F^2<\infty$. The population mean function is
\[
    \mu(t)=\E[F(t)],
\]
where the pointwise expression is interpreted when evaluation is meaningful.
As an element of $L^2$, the sample mean is
\[
    \widehat\mu(t)=\frac1n\sum_{i=1}^nf_i(t).
\]

\begin{proposition}[Mean as a Fréchet mean]
The mean $\mu=\E F$ uniquely minimizes
\[
    g\longmapsto\E\norm{F-g}_{L^2}^2.
\]
Similarly, $\widehat\mu$ minimizes
$n^{-1}\sum_i\norm{f_i-g}_{L^2}^2$.
\end{proposition}

\begin{proof}
The Hilbert-space identity
\[
 \E\norm{F-g}^2
 =\E\norm{F-\mu}^2+\norm{\mu-g}^2
\]
follows because $\E(F-\mu)=0$ makes the cross term vanish.
\end{proof}

This variational definition generalizes to nonlinear spaces: replace the
Hilbert norm by an appropriate metric and minimize the sum of squared
distances. Uniqueness is no longer automatic.

\subsection{Covariance Function}

\begin{definition}[Covariance function]
For a square-integrable process $F$,
\[
    K(s,t)=\Cov\{F(s),F(t)\}
    =\E[(F(s)-\mu(s))(F(t)-\mu(t))].
\]
\end{definition}

The covariance is symmetric and nonnegative definite:
\[
 K(s,t)=K(t,s),\qquad
 \sum_{i,j=1}^ma_ia_jK(t_i,t_j)\geq0.
\]
Its empirical estimate from densely observed curves is
\[
    \widehat K(s,t)=\frac1{n-1}\sum_{i=1}^n
    \{f_i(s)-\widehat\mu(s)\}\{f_i(t)-\widehat\mu(t)\}.
\]

The diagonal $K(t,t)$ is pointwise variance, while off-diagonal values
describe association across locations. Unlike a finite covariance matrix,
$K$ is a function on a two-dimensional domain.

\section{Linear Operators on Function Spaces}

\begin{definition}[Bounded linear operator]
For normed spaces $V,W$, a linear map $T:V\to W$ is bounded if there is
$C<\infty$ such that
\[
    \norm{Tf}_W\leq C\norm f_V\qquad\text{for every }f\in V.
\]
Its operator norm is
$\norm T_{\mathrm{op}}=\sup_{\norm f\leq1}\norm{Tf}$.
\end{definition}

For linear maps between normed spaces, boundedness is equivalent to
continuity. In finite dimensions every linear map is bounded; this is false
in infinite dimensions. For example, differentiation is not bounded from
$C^1([0,1])$ equipped only with the sup norm into $C([0,1])$.

\subsection{Integral Operators}

Given a kernel $k:[0,1]^2\to\R$, define
\begin{equation}
    (Tf)(t)=\int_0^1k(s,t)f(s)\dd s.
    \label{eq:integral-operator}
\end{equation}
This map is linear. If $k\in L^2([0,1]^2)$, then $T$ is a bounded
Hilbert--Schmidt operator on $L^2$ and
\[
    \norm T_{\mathrm{op}}\leq
    \left\{\int_0^1\int_0^1k(s,t)^2\dd s\dd t\right\}^{1/2}.
\]

\begin{definition}[Adjoint]
For Hilbert spaces $H_1,H_2$ and bounded $T:H_1\to H_2$, the adjoint
$T^*:H_2\to H_1$ is the unique operator satisfying
\[
    \inner{Tf}{g}_{H_2}=\inner f{T^*g}_{H_1}
\]
for all $f,g$.
\end{definition}

For matrices, the adjoint is the transpose. For the integral operator
\eqref{eq:integral-operator},
\[
    (T^*g)(s)=\int_0^1k(s,t)g(t)\dd t.
\]
Thus $T$ is self-adjoint when $k(s,t)=k(t,s)$ almost everywhere.

\begin{definition}[Positive operator]
A self-adjoint operator $T$ is positive semidefinite if
$\inner{Tf}{f}\geq0$ for every $f$.
\end{definition}

\subsection{Covariance Operator}

The covariance function induces
\begin{equation}
    (\mathcal Kf)(t)=\int_0^1K(s,t)f(s)\dd s.
    \label{eq:covariance-operator}
\end{equation}
Equivalently, in Hilbert notation,
\[
    \mathcal Kf
    =\E[\inner{F-\mu}{f}(F-\mu)].
\]

\begin{proposition}
The covariance operator is self-adjoint and positive semidefinite. If
$\E\norm{F-\mu}^2<\infty$, it is trace class and
\[
    \tr(\mathcal K)=\E\norm{F-\mu}^2
                  =\int_0^1K(t,t)\dd t.
\]
\end{proposition}

For any $f$,
\[
 \inner{\mathcal Kf}{f}
 =\E[\inner{F-\mu}{f}^2]\geq0.
\]
This is the functional analogue of $a^\mathsf T\Sigma a\geq0$.

\section{Principal Component Analysis}

For centered vectors $X\in\R^p$ with covariance $\Sigma$, the first principal
direction solves
\[
    v_1=\argmax_{\norm v=1}\Var(v^\mathsf TX)
       =\argmax_{\norm v=1}v^\mathsf T\Sigma v.
\]

\begin{theorem}[Rayleigh quotient]
For symmetric $\Sigma$, the maximum of
$v^\mathsf T\Sigma v$ over $\norm v=1$ is the largest eigenvalue
$\lambda_1$, attained by a corresponding eigenvector. Subsequent principal
directions are eigenvectors obtained under orthogonality constraints.
\end{theorem}

For observations $x_1,\ldots,x_n$, the rank-$q$ PCA subspace also minimizes
reconstruction error:
\[
 \min_{U^\mathsf TU=I_q}\sum_{i=1}^n
 \norm{x_i-UU^\mathsf Tx_i}^2.
\]
If the centered data matrix has SVD $X=UDV^\mathsf T$, loadings are columns of
$V$, scores are columns of $UD$, and explained variances are proportional to
$d_j^2$.

\section{Functional Principal Component Analysis}

Functional PCA (FPCA) replaces the covariance matrix by the covariance
operator. A principal direction $\phi$ is a unit-norm function maximizing
\[
    \Var\{\inner{F-\mu}{\phi}\}
    =\inner{\mathcal K\phi}{\phi}.
\]
The solution satisfies the eigenequation
\begin{equation}
    \mathcal K\phi_j=\lambda_j\phi_j,
    \qquad
    \int_0^1K(s,t)\phi_j(s)\dd s
    =\lambda_j\phi_j(t).
    \label{eq:fpca-eigenproblem}
\end{equation}

\begin{theorem}[Spectral decomposition of covariance]
Let $\mathcal K$ be a compact, self-adjoint, positive covariance operator.
Then it has orthonormal eigenfunctions $\{\phi_j\}$ with nonnegative
eigenvalues $\lambda_1\geq\lambda_2\geq\cdots\to0$, and
\[
    K(s,t)=\sum_{j=1}^\infty\lambda_j\phi_j(s)\phi_j(t)
\]
in the appropriate $L^2$ sense.
\end{theorem}

\subsection{Karhunen--Loève Expansion}

Define the principal component scores
\[
    \xi_{ij}=\inner{f_i-\mu}{\phi_j}
    =\int_0^1\{f_i(t)-\mu(t)\}\phi_j(t)\dd t.
\]

\begin{theorem}[Karhunen--Loève]
Under the preceding second-moment assumptions,
\[
    F(t)=\mu(t)+\sum_{j=1}^\infty\xi_j\phi_j(t)
\]
in mean square, where
\[
    \E\xi_j=0,\qquad
    \Cov(\xi_j,\xi_k)=\lambda_j\ind\{j=k\}.
\]
\end{theorem}

Truncation gives the model
\begin{equation}
    f_i(t)\approx\mu(t)+\sum_{j=1}^q\xi_{ij}\phi_j(t).
    \label{eq:truncated-fpca}
\end{equation}
The fraction of integrated variance explained by the first $q$ components is
\[
    \frac{\sum_{j=1}^q\lambda_j}
         {\sum_{j=1}^\infty\lambda_j}.
\]

\begin{proposition}[Optimal reconstruction]
Among all $q$-dimensional linear subspaces of $L^2$, the span of
$\phi_1,\ldots,\phi_q$ minimizes expected squared reconstruction error, which
equals $\sum_{j>q}\lambda_j$.
\end{proposition}

\subsection{Computation from a Common Grid}

Suppose every centered curve is sampled on a common grid
$t_1,\ldots,t_m$. Stack curve vectors as rows of
$X_c\in\R^{n\times m}$. Ordinary PCA of $X_c$ approximates FPCA only after
accounting for numerical integration. With quadrature weights
$W=\diag(w_1,\ldots,w_m)$, the discrete inner product is
\[
    \inner f g\approx f^\mathsf TWg.
\]
An SVD of $X_cW^{1/2}$ gives eigenfunctions normalized in the weighted inner
product. Ignoring unequal grid spacing changes the geometry.

Alternatively, if $f_i(t)=b(t)^\mathsf Tc_i$ and
$G=\int b(t)b(t)^\mathsf T\dd t$, then functional orthonormality is measured
by $G$, not the Euclidean coefficient norm. FPCA becomes a generalized
eigenproblem or an ordinary eigenproblem after whitening by $G^{1/2}$.

\begin{example}[Interpreting a mode]
Plot
\[
    \widehat\mu(t)\pm c\sqrt{\widehat\lambda_j}\,
    \widehat\phi_j(t)
\]
for a small constant $c$, often one or two. These curves visualize the
functional change associated with positive and negative scores. The sign of
an eigenfunction is arbitrary, so interpretation must be invariant to a
simultaneous sign change of $\phi_j$ and its scores.
\end{example}

\subsection{Practical Qualifications}

\begin{itemize}
    \item Measurement error inflates the empirical covariance diagonal.
    Smooth the covariance surface or use a model that separates signal and
    noise rather than treating the diagonal spike as functional variation.
    \item Dense common grids permit direct covariance estimation. Sparse,
    irregular designs require pooled smoothing and conditional score
    estimation.
    \item Eigenfunctions with nearly tied eigenvalues are individually
    unstable; their joint eigenspace can still be stable.
    \item Standard FPCA measures amplitude variation under $L^2$. Timing
    variation can dominate the components unless curves are registered.
\end{itemize}

\end{document}
